\documentclass[11pt]{article}
\usepackage[T1]{fontenc}
\usepackage{textcomp}
\usepackage[margin=0.75in]{geometry}
\usepackage{amsmath,amssymb,amsthm}
\usepackage{array,booktabs}
\usepackage{hyperref}
\setlength{\parindent}{0pt}
\newtheorem{theorem}{Theorem}
\theoremstyle{definition}
\newtheorem{definition}{Definition}
\title{Flow Proof Artifact: circle-radii-equal}
\author{Generated by \texttt{flow doc proof}}
\date{Flow Proof Book}
\begin{document}
\maketitle
All radii of a circle from the centre to the circumference are equal.\\[0.5em]
\textbf{Source.} euclid — Elements, Book I, Definition 15\\[0.5em]
\bigskip
\noindent{\large \textbf{Definition 1.} \textit{All radii of a circle from the centre to the circumference are equal} \hfill the Euclidean plane \cdot circle \cdot radii from the centre are equal}\par
\smallskip\noindent\textit{Coordinate: the Euclidean plane \cdot circle \cdot radii from the centre are equal}\par
\smallskip\noindent\textit{Source: euclid — Elements, Book I, Definition 15}\par
\medskip
\noindent\textbf{Goal.} All radii of a circle from the centre to the circumference are equal\par
\noindent$\displaystyle OA = OB$\par
\medskip
\noindent
\begin{tabular}{@{} >{\raggedright\arraybackslash}p{0.06\textwidth} >{\raggedright\arraybackslash}p{0.47\textwidth} >{\raggedleft\arraybackslash}p{0.41\textwidth} @{}}
\textbf{\#} & \textbf{Proof} & \textbf{Mathematics} \\
\midrule
\textbf{1.} & We stipulate radii from the centre are equal for circles in the Euclidean plane — this is a definition, not a derived fact. & \\[0.45em]
\textbf{2.} & This follows directly from the definition: radius OA equals radius OB. Hence proven. & $\displaystyle OA = OB$ \\[0.45em]
\bottomrule
\end{tabular}
\label{thm:Geometry-circle-radii_from_the_centre_are_equal}
\par\medskip\hrule
\end{document}
